\documentclass[11pt]{article}
\usepackage{mathblatt}
\newcommand{\bereich}[1]{\par\addvspace{14pt}\noindent{\large\bfseries #1}\par\nobreak\vspace{4pt}}
\newcommand{\bl}[1]{\par\noindent #1\par\vspace{3pt}}
\newcommand{\blrand}[1]{\par\noindent{\color{mbgrau}#1}\par\vspace{3pt}}
\newcommand{\blhier}[1]{\par\noindent\llap{$\blacktriangleright$\hspace{0.5em}}\textbf{#1}\par\vspace{3pt}}
\newcommand{\stufe}[1]{\par\noindent\hspace*{1.5em}{\small #1}\par\vspace{2pt}}
\newcommand{\pkt}{\textperiodcentered{}}
\begin{document}

\pfheftstil
\pflaufkopf[Abitur GK \pkt{} Prüfung]{Integral}{AA150}\pfseitevon
\pfheftkopf{Integral}{}
\bl{Integral berechnen}
\blhier{Flächeninhalt}
\stufe{Fläche gesucht}
\stufe{Fläche im Sachzusammenhang gesucht}
\stufe{Lösungsweg erklären}
\bl{Integral am Graphen}

\par\addvspace{14pt}\noindent Vorher: Ableitung + Tangente\par
\par\noindent Weiter: Punkte, Flächen, Körper\par
\end{document}